\documentclass[a4paper,fleqn,doubleblind]{cas-dc}
\usepackage[T1]{fontenc}
\usepackage[utf8]{inputenc}
\usepackage{charter}
\renewcommand{\sfdefault}{bch}
\usepackage{textcomp}
\usepackage{amsmath,amssymb}
\usepackage[numbers,sort&compress]{natbib}
\usepackage{graphicx,booktabs,array,enumitem}
\usepackage{xurl}
\urlstyle{same}
\usepackage[switch]{lineno}
\usepackage[expansion=false]{microtype}
\def\tsc#1{\csdef{#1}{\textsc{\lowercase{#1}}\xspace}}
\renewcommand{\linenumberfont}{\normalfont\tiny\color{gray}}
\setlength{\bibsep}{1pt}
\renewcommand{\figurename}{Fig.}
\begin{document}
\let\WriteBookmarks\relax
\let\printorcid\relax
\def\floatpagepagefraction{1}
\def\textpagefraction{.001}
\shorttitle{Self-timed end-of-life alarms}
\shortauthors{Anonymized manuscript}
\title [mode = title]{Knowing when to warn: Self-timed end-of-life alarms under an asymmetric tolerance window, their maintenance cost, and three pre-specified confirmations}
\begin{abstract}
Remaining useful life (RUL) models are usually scored at fixed times, yet a replacement decision is taken once. On PHM08 engine data, none of eleven models placed more than 33.5\% of half-life predictions within an asymmetric tolerance window of \textminus{}13 to +2 cycles. We propose a self-timed warning model: a horizon-focused gradient-boosted regressor with a window-centred, debounced trigger, tuned by nested cross-validation, that raises one scored alarm per unit. At a warning level of 20 cycles, 91.7\% of PHM08 and 90.2\% of 927 simulated engine alarms fell inside the window. Controls show what this measures: a constant estimate at the same alarm times scored 90.8\%, whereas a health-indicator threshold alarm scored 53.4\%; the regressor's value lies in timing the alarm. In a renewal--reward model, window-tuned alarms were cheaper than optimal age replacement in 94 of 96 lead-time and cost settings on engines and batteries, and had lower mean regret than alarms tuned directly to cost (5.5\% against 10.9\%). Three pre-specified confirmations (protocols hashed, not registered) delimit the method: 97.6\% on 123 lithium-ion cells, 68.7\% on N-CMAPSS and 0.0\% on 17 rolling bearings.
\end{abstract}
\begin{keywords}
Remaining useful life \sep
Predictive maintenance \sep
Early warning \sep
Alarm timing \sep
Renewal--reward model \sep
Turbofan engine \sep
Lithium-ion battery
\end{keywords}
\maketitle
\linenumbers
\section{Introduction}\label{sec:1}

Data-driven prognostics estimates the remaining useful life (RUL) of a component from its condition-monitoring history, so that maintenance can be scheduled before failure \cite{lei2018,si2011}. Most studies on turbofan-engine benchmarks report an average error over prediction times chosen by the evaluator: the last observed cycle of truncated test trajectories, all cycles of each trajectory, or fixed fractions of life \cite{ramasso2014}. A maintenance decision, however, is taken at one moment. What matters to the operator is that the estimate available at that moment is accurate enough to act on, and that it is not late, because a late estimate lets an engine run into failure.

This study started from a concrete requirement of that kind. A remaining-life estimate was to be counted as correct only if its error, predicted minus true cycles left, lay within \textminus{}13 to +2 cycles: up to 13 cycles early was tolerable, more than 2 cycles late was not. The early limit coincides with the early-prediction constant of the PHM08 challenge score \cite{saxena2008}; the late limit is much stricter than that score's late constant of 10. We first asked how often existing models meet this requirement when the prediction is made at a fixed point, half of each engine's life. The answer was sobering. Across eleven models, from gradient-boosted trees and recurrent and convolutional networks to a Transformer, a temporal convolutional network with generative data augmentation, a physics-informed Bayesian degradation model and a stacked hybrid of most of them, the best placed 33.5\% of PHM08 predictions inside the window (Section 6.1). At half of life the sensors show almost no wear: the change of a pooled health signal from its early-life level (median 0.23) is comparable to its cycle-to-cycle noise (0.21), whereas 20 cycles before failure it is 1.08. Locating the end of life to within 15 cycles from such signals is not realistic.

The requirement becomes attainable if the model is allowed to choose when to speak. In service, an engine is monitored at every cycle, and the decision that needs a precise estimate is the final one: when to remove the engine or to schedule the shop visit. We therefore reformulate the task as a first-warning problem. A policy observes each engine cycle by cycle, stays silent while the engine looks healthy, and issues a single warning with a remaining-life estimate when it judges the end to be near. The warning is scored at the moment it is issued; an engine that fails without a warning counts as late. This is close in spirit to the prognostic horizon of Saxena et al. \cite{saxena2010}, the time from which predictions stay within an accuracy band, and to alarm-based maintenance planning \cite{depater2022}, but it evaluates one actionable estimate per engine instead of a curve of predictions. Two questions follow, and this paper answers both: what such a score actually measures, and whether tuning for it gives good maintenance decisions. We find that the score measures \textit{when} the alarm is raised far more than the number attached to it, and that alarms tuned for the window are a robust choice when priced by the long-run cost of replacements and failures.

The contributions of this paper are as follows.

\begin{itemize}[leftmargin=1.2em,itemsep=2pt]
\item[(i)] A \textbf{self-timed warning (STW) model} (Fig. 1): a condition-normalised, causal feature encoder; a gradient-boosted regressor with a capped remaining-life target; and a window-centred, debounced trigger. All four decision settings are chosen by nested cross-validation on training units, so test units never influence the model.
\item[(ii)] \textbf{What the score measures.} A constant-estimate control and a health-indicator threshold alarm show that the in-window rate is a property of the alarm time: replacing the model's estimate by a constant did not lower the rate on any of seven datasets, while timing the alarm on a classical health indicator lowered it to about half. A window-centring result (Proposition 1) explains the role of the shift and why the attainable rate is set by how concentrated the alarm time is relative to failure.
\item[(iii)] \textbf{A cost evaluation of the alarms.} In a renewal--reward model with lead times and failure costs, window-tuned alarms are compared with alarms tuned directly to the cost, with a cost-tuned health-indicator alarm and with optimal age replacement; the window turns out to be a robust surrogate for the cost.
\item[(iv)] A \textbf{leakage-safe first-warning evaluation} on five simulated engine benchmarks (90.2\% of 927 first warnings in the window at a warning level of 20 cycles), with an identical-protocol comparison against re-implementations of the three winning PHM08 challenge approaches and deep networks given the same inputs, and a component ablation.
\item[(v)] \textbf{Three pre-specified confirmations}, each with the protocol fixed and hashed before the first run: 97.6\% on measured run-to-failure data of 123 lithium-ion cells \cite{severson2019}, with all four hypotheses supported; 68.7\% on N-CMAPSS \cite{ariaschao2021}, with three of four unsupported; and 0.0\% on 17 rolling bearings \cite{nectoux2012}, with no hypothesis supported. Together they show that the method works on some systems and not on others.
\end{itemize}

Section 2 reviews related work; Section 3 defines the task; Section 4 describes the model and the window-centring result; Section 5 gives the experimental design; Section 6 reports results; Section 7 discusses them; Section 8 sets out future work on test and maintenance policies and on automatic fault detection and diagnosis; and Section 9 concludes.

\section{Literature review and contribution}\label{sec:2}

\subsection{Remaining useful life prediction for turbofan engines}

The C-MAPSS simulations and the PHM08 challenge data \cite{saxena2008} are the standard benchmarks for engine prognostics. The three winning PHM08 challenge approaches, published in the conference proceedings, were a similarity-based method that matches the test trajectory against training trajectories \cite{wang2008}, a recurrent neural network \cite{heimes2008} and an ensemble of neural networks smoothed by a Kalman filter \cite{peel2008}. Later work moved to deep sequence models, long short-term memory networks \cite{zheng2017} and convolutional networks \cite{li2018}, usually trained on a remaining-life target capped at about 125 cycles. Ramasso and Saxena \cite{ramasso2014} catalogued the evaluation practices on these data and showed how strongly reported scores depend on the choice of test cut-off points and metrics. N-CMAPSS \cite{ariaschao2021} adds realistic flight profiles to the simulation. On tabular inputs, tree ensembles remain strong competitors to deep networks \cite{grinsztajn2022}, which motivated our choice of gradient boosting \cite{friedman2001}. Lithium-ion batteries pose the same question in a different physical setting: Severson et al. \cite{severson2019} cycled commercial cells to failure under fast-charging protocols and predicted their cycle life from the first cycles; their data are one of the few public run-to-failure sets with real measurements and more than a hundred units.

\subsection{Evaluation at the time of decision}

Saxena et al. \cite{saxena2010} proposed metrics that follow a prediction over time, among them the prognostic horizon and $\alpha$--$\lambda$ accuracy, where the tolerated error shrinks as failure approaches. Asymmetric losses that penalise late predictions more than early ones are common \cite{rengasamy2020}, and the PHM08 score itself is asymmetric. Maintenance-oriented studies link prognostics to decisions: alarm-based scheduling for aircraft engines \cite{depater2022}, dynamic maintenance frameworks driven by predicted failure probabilities \cite{nguyen2019,mitici2023}, reinforcement learning on probabilistic prognostics \cite{lee2023} and decision-based metrics for prognostic algorithms \cite{kamariotis2024}. Degradation-signal models with Bayesian updating \cite{gebraeel2005} produce a residual-life distribution that can be acted on when needed. The classical reference for any condition-based policy is age replacement, which uses no condition information and replaces a unit at the age that minimises the long-run cost rate \cite{barlow1960}; we use it as the baseline of our cost analysis. Our contribution sits between these lines: we keep a standard regression model but let a tuned decision rule choose the single moment at which its estimate is issued and judged. Unlike alarm-based scheduling, which acts when a probabilistic RUL estimate crosses a threshold \cite{depater2022}, and decision-oriented metrics, which score a prognostic algorithm by the cost of the decisions it induces \cite{kamariotis2024}, we ask what a tolerance-window score rewards, show that it rewards the alarm time rather than the estimate, and measure how well window-tuned alarms perform when priced by cost.

\subsection{Leakage and model selection}

Model selection on the data used for evaluation biases error estimates \cite{varma2006}, and leakage between training and test data is a leading cause of irreproducible results in machine-learning-based science \cite{kapoor2023}. Run-to-failure benchmarks invite two specific leaks: splitting by rows rather than by engines, and evaluating at a cut point derived from the engine's total life, which the model could exploit. Our design avoids both, and Section 5.5 describes the checks. Pre-specifying an analysis and fixing it before the data are seen is the standard remedy against selective reporting in confirmatory studies \cite{nosek2018}; we apply it to three confirmation datasets.

\section{Problem formulation}\label{sec:3}

An engine \textit{i} is observed at flight cycles \textit{t} = 1, \ldots{}, \textit{L}\textsubscript{i}, where \textit{L}\textsubscript{i} is its (unknown) life. At cycle \textit{t} the true remaining life is \textit{r}\textsubscript{t} = \textit{L}\textsubscript{i} \textminus{} \textit{t}. A warning policy maps the history observed up to \textit{t} to either \textit{silent} or \textit{warn}, together with an estimate $\hat{y}$\textsubscript{t}. The first-warning time is

\begin{equation}
t^{*} = \min\{\, t < L_i : \text{the policy warns at } t \,\},
\end{equation}

and the error of the warning is \textit{e} = $\hat{y}$\textsubscript{t*} \textminus{} \textit{r}\textsubscript{t*}. A warning is \textbf{correct} if \textminus{}13 $\le$ \textit{e} $\le$ +2, \textbf{too early} if \textit{e} < \textminus{}13 and \textbf{too late} if \textit{e} > +2. An engine that reaches failure without a warning is \textbf{missed} and counted with the late outcomes, since it is the most dangerous case. The primary metric is the share of engines whose first warning is correct (the in-window rate). We also report the late-or-missed and early rates, the lead time \textit{r}\textsubscript{t*}, and the share of warnings for which the early limit can bind at all. A warning issued with fewer than 13 cycles left can be too early only if its estimate is negative, which practically never happens. If a policy warned only in the last few cycles, the early limit would be inactive and the task much easier. We therefore report, for every setting, the share of warnings issued with at least 13 cycles left. We write ``engine'' for the monitored unit; on the battery and bearing data a unit is a cell or a bearing, and a cycle of the model is a block of consecutive charge--discharge cycles or vibration snapshots (Section 5.4).

The warning level \textit{H} expresses when the operator wants to be warned: the policy warns when it believes that about \textit{H} cycles are left. Small \textit{H} gives accurate but short-notice warnings; large \textit{H} gives more notice but less precision. We report \textit{H} $\in$ \{10, 15, 20\} on all datasets and also \textit{H} = 30 and 40 on PHM08 (and on all engine datasets in the secondary analyses of Section 5.6). \textit{H} = 20 is the primary level because it is the smallest of these at which the early limit can bind for more than 90\% of PHM08 warnings (Sections 5.2 and 6.3).

Two remarks frame the evaluation. First, a warning is correct exactly when the true remaining life at the alarm lies between $\hat{y}$\textsubscript{t*} \textminus{} 2 and $\hat{y}$\textsubscript{t*} + 13. The window therefore scores the alarm time and the estimate jointly, and Section 6.3 shows that, for our model, the alarm time carries almost all of the information. Second, the window is an operational acceptance criterion, not a cost. We therefore also price every alarm policy in a renewal--reward model with replacement lead times and failure costs (Sections 5.6 and 6.6).

\section{The self-timed warning model}\label{sec:4}

Fig. 1 shows the architecture. The online part runs at every cycle of an engine in service; the offline part fixes its settings using training engines only.

\begin{figure*}[!tbp]\centering
\includegraphics[width=\textwidth]{img/Fig1_STW_architecture.png}
\caption{Architecture of the self-timed warning (STW) model. Online (top): regime-normalised sensors are encoded causally and mapped to an estimate $\hat{y}$\textsubscript{t}, which is smoothed and shifted to $\tilde{y}$\textsubscript{t}; a warning is issued the first time $\tilde{y}$ stays at or below \textit{H} for \textit{k} cycles. Offline (bottom): $\theta$ = (\textit{c}, \textit{m}, \textit{k}, \textit{s}) is chosen on inner out-of-fold predictions of the training engines and frozen; test engines are scored once. Inset: the scoring window for \textit{e} = $\tilde{y}$\textsubscript{t*} \textminus{} \textit{r}\textsubscript{t*}.}
\end{figure*}

\subsection{Condition normalisation}

Engine sensor readings depend strongly on the operating condition (altitude, Mach number, throttle). Each sensor \textit{j} is z-scored within its operating regime \textit{a}, with means and standard deviations estimated from training engines only:

\begin{equation}
z_{t,j} = \bigl( x_{t,j} - \mu_j(a_t) \bigr) / \sigma_j(a_t).
\end{equation}

PHM08, FD002 and FD004 have six regimes, FD001 and FD003 one, and N-CMAPSS uses the three flight classes. A sensor is dropped if it is constant within any regime, which leaves 15 sensors on PHM08, FD001 and FD002 and 16 on FD003 and FD004. On N-CMAPSS we used, as in the pre-specified protocol, the six temperature and shaft-speed channels T24, T30, T48, T50, Nf and Nc; an exploratory re-run with all 14 channels is reported in Section 6.7. On the battery data, the seven per-cycle quantities (discharge and charge capacity, internal resistance, average, minimum and maximum temperature, and charge time) are z-scored as one regime, as fixed in the protocol. The protocol justified this by stating that all cells share the same test conditions apart from the charging protocol; in fact the three batches also differ in the rest periods of the cycling schedule \cite{severson2019}. The model is given neither.

\subsection{Causal feature encoder}

At cycle \textit{t} the encoder uses only cycles 1, \ldots{}, \textit{t}. It computes an early-life baseline \textbf{b}\textsubscript{t}, the mean of the first min(20, \textit{t}) cycles, which captures manufacturing variation and initial wear; exponentially weighted means with spans 10 and 30 of the baseline-removed signals \textbf{z} \textminus{} \textbf{b}; least-squares slopes over the last 15 and 40 cycles (smoothed with span 5), which capture the rate of degradation; and the age \textit{t}. This gives 5\textit{p} + 1 features for \textit{p} sensors. Section 5.5 verifies causality: deleting all data after cycle \textit{t} does not change the prediction at \textit{t}.

\subsection{Horizon-focused regressor}

A histogram gradient-boosted regressor \cite{pedregosa2011} with absolute-error loss is trained on all cycles of the training engines with the target

\begin{equation}
y_t = \min( r_t, c ), \qquad c \in \{40, 60, 90\}.
\end{equation}

The conventional piecewise-linear target caps remaining life at about 125 cycles because early degradation is unobservable \cite{heimes2008}. We allow lower caps, so that above \textit{c} the model only has to say ``far from failure'' and its loss is spent on the last \textit{c} cycles, where the warning is issued; a cap also flattens the prediction during the healthy phase. The cap is chosen per split by the inner cross-validation (Section 4.5). The ablation in Section 6.5 shows that having a cap matters, whereas its exact value matters little: the conventional cap of 125 performed almost as well as the selected caps of 40 or 60.

\subsection{Window-centred, debounced trigger}

The raw estimate is smoothed causally and shifted:

\begin{equation}
\tilde{y}_t = \mathrm{EWM}_m(\hat{y})_t + s,
\end{equation}

where EWM\textsubscript{m} is an exponentially weighted mean with span \textit{m} $\in$ \{1, 3, 6\} (\textit{m} = 1 means no smoothing) and \textit{s} $\in$ [\textminus{}12, +8] in steps of 0.5. The shift centres the estimates in the tolerance window: the window [\textminus{}13, +2] is centred at \textminus{}5.5, and a negative shift trades the narrow late side against the wide early side (Section 4.6 makes this precise). A warning is issued at the first cycle at which $\tilde{y}$ has been at or below \textit{H} for \textit{k} $\in$ \{1, 2, 3, 5\} consecutive cycles:

\begin{equation}
t^{*} = \min\{\, t : \tilde{y}_{\tau} \le H \ \text{for all}\ \tau = t-k+1, \ldots, t \,\},
\end{equation}

and the reported estimate is $\tilde{y}$\textsubscript{t*}. The debounce \textit{k} is meant to suppress warnings triggered by a single noisy cycle; the ablation shows that its effect, like that of the smoothing, is small.

\subsection{Nested selection of the decision settings}

The settings $\theta$ = (\textit{c}, \textit{m}, \textit{k}, \textit{s}) are chosen on training engines only. The training engines of an outer split are divided into five inner folds (seed 100 + outer fold). For each inner fold and each cap, the regressor is fitted on the other inner folds, which also supply the normalisation statistics, and predicts every cycle of the held-out inner engines. This gives out-of-fold predictions for all training engines. For each warning level \textit{H}, the grid (the caps with \textit{c} $\ge$ \textit{H} + 20 $\times$ 3 spans $\times$ 4 debounces $\times$ 41 shifts) is searched for the setting that maximises the number of training engines whose first warning is correct; ties are broken towards smaller |\textit{s}|. The regressor is then refitted on all training engines with the selected cap, $\theta$ is frozen, and each test engine is processed once. The search is vectorised over \textit{s}. Supplementary Table S8 lists all settings, and Supplementary Section S14 lists the scripts that produce every number.

\subsection{Why a shift, and what limits the in-window rate}

A simple model explains the role of the shift and gives a benchmark for the attainable rate. Let $\varepsilon$ = EWM\textsubscript{m}($\hat{y}$)\textsubscript{t*} \textminus{} \textit{r}\textsubscript{t*} be the error of the unshifted (smoothed) estimate at the moment of the first warning, so that the scored error is \textit{e} = $\varepsilon$ + \textit{s}, and write the window as [\textit{a}, \textit{b}] = [\textminus{}13, +2] with width \textit{w} = \textit{b} \textminus{} \textit{a} = 15.

\textbf{Proposition 1.} Suppose that, over the range of shifts considered, $\varepsilon$ has a continuous distribution \textit{F} that does not depend on \textit{s} (assumption A1). Then the probability of a correct first warning is

\begin{equation}
P(s) = F(b - s) - F(a - s).
\end{equation}

(i) For every \textit{s}, \textit{P}(\textit{s}) $\le$ \textit{C}\textsubscript{w}(\textit{F}) = sup\textsubscript{x} [\textit{F}(\textit{x} + \textit{w}) \textminus{} \textit{F}(\textit{x})], the largest probability that $\varepsilon$ falls in some interval of width \textit{w} (the \textit{concentration} of the error). (ii) If \textit{F} has a density that is unimodal and symmetric about its median \textit{m}, \textit{P}(\textit{s}) is maximised at

\begin{equation}
s^{*} = (a + b)/2 - m = -5.5 - m,
\end{equation}

and the maximum is \textit{C}\textsubscript{w}(\textit{F}) = Pr(|$\varepsilon$ \textminus{} \textit{m}| $\le$ \textit{w}/2) = Pr(|$\varepsilon$ \textminus{} \textit{m}| $\le$ 7.5). The proof is given in Supplementary Section S2.

Two consequences guide the design and the interpretation. First, an unbiased estimator (\textit{m} $\approx$ 0) should be shifted by about \textminus{}5.5 cycles: without the shift every error above +2 cycles is late, a large part of the upper half of the error distribution, while most of the wide early side goes unused. This is why removing the shift is the most damaging single-component ablation (Section 6.5). Second, once the shift is right, the in-window rate is governed by the concentration of the error \textit{at the warning}, not by the average error over life. This is what the horizon-focused target is for: capping the target at \textit{c} spends the regressor's capacity on the last \textit{c} cycles and tightens $\varepsilon$ near \textit{H}. Assumption A1 is an approximation, because changing \textit{s} moves the warning to a slightly different cycle; Section 6.8 checks the Proposition against all datasets, and Supplementary Section S2 gives the proof.

\section{Experimental design}\label{sec:5}

\subsection{Data}

\begin{table*}[!tbp]
\footnotesize\setlength{\tabcolsep}{4pt}
\caption{Datasets. Life = number of cycles to failure of the run-to-failure units.}
\begin{tabular}{@{}>{\raggedright\arraybackslash}p{\dimexpr0.1662\linewidth-2\tabcolsep\relax}>{\raggedright\arraybackslash}p{\dimexpr0.1053\linewidth-2\tabcolsep\relax}>{\raggedright\arraybackslash}p{\dimexpr0.1108\linewidth-2\tabcolsep\relax}>{\raggedright\arraybackslash}p{\dimexpr0.2105\linewidth-2\tabcolsep\relax}>{\raggedright\arraybackslash}p{\dimexpr0.1329\linewidth-2\tabcolsep\relax}>{\raggedright\arraybackslash}p{\dimexpr0.2743\linewidth-2\tabcolsep\relax}@{}}
\toprule
Dataset & Units & Rows & Life (min--max; median) & Regimes & Role \\
\midrule
PHM08 & 218 & 45,918 & 128--357; 209.5 & 6 & development \\
FD001 & 100 & 20,631 & 128--362; 199 & 1 & additional evaluation \\
FD002 & 260 & 53,759 & 128--378; 199 & 6 & additional evaluation \\
FD003 & 100 & 24,720 & 145--525; 220.5 & 1 & additional evaluation \\
FD004 & 249 & 61,249 & 128--543; 234 & 6 & additional evaluation \\
N-CMAPSS & 99 & 7,473 & 48--100; 75 & 3 (flight class) & pre-specified confirmation 1 (simulation) \\
Battery \cite{severson2019} & 123 cells & 95,276 & 148--1935; 731 & 1 & pre-specified confirmation 2 (laboratory measurements) \\
Bearings \cite{nectoux2012} & 17 bearings & 24,889 & 230--2803; 1428 & 3 (load/speed) & pre-specified confirmation 3 (laboratory measurements) \\
\bottomrule
\end{tabular}\par\vspace{3pt}{\raggedright\scriptsize N-CMAPSS: one row per engine and flight cycle (means of the 1 Hz signals), nine subsets; 60 development and 39 test engines in the official split. Battery: one row per cell and charge--discharge cycle; life = official cycle life (first cycle below 80\% of the nominal 1.1 Ah capacity); 123 of the 140 cell records (135 physical cells, because five records continue five others) after the exclusions of Supplementary Section S5. Bearings: one row per vibration snapshot (0.1 s every 10 s); life = number of snapshots to failure.\par}
\end{table*}

We use the 218 run-to-failure training engines of PHM08 and the training sets of C-MAPSS FD001--FD004 \cite{saxena2008}. The truncated official test sets are not used, because the first-warning task needs complete trajectories. For N-CMAPSS \cite{ariaschao2021}, the 1 Hz flight data were reduced to one row per engine and flight cycle (the mean of each signal over the flight). All 99 engines of the nine subsets present in the archive we used (DS08d was absent) were used. The battery data are laboratory measurements of commercial lithium iron phosphate/graphite cells (A123 APR18650M1A, 1.1 Ah) cycled to failure under fast-charging protocols in three batches (\cite{severson2019}; 66 distinct protocols among the cells used); we used the per-cycle summary of each cell. The bearing data are the 17 complete run-to-failure records of the PRONOSTIA platform \cite{nectoux2012}, with horizontal and vertical acceleration sampled at 25.6 kHz in 0.1-s snapshots every 10 s.

\subsection{Development history and splits}

The model was developed on PHM08 after the fixed-time results of the age-only, gradient-boosted, recurrent, convolutional, hidden-Markov and similarity models were known; the other fixed-time reference models of Section 6.1 (Transformer, Bayesian, temporal convolutional and stacked models) were run afterwards, in an attempt to improve fixed-time prediction. The warning levels and grids were fixed on PHM08, and C-MAPSS FD001--FD004 were then run once with unchanged code; we call them additional evaluation datasets, although they were not pre-specified. The primary level \textit{H} = 20 was designated after these runs, by the criterion of Section 3 computed on PHM08 (the early limit could bind for 95.4\% of warnings at \textit{H} = 20 and 65.6\% at \textit{H} = 15); the C-MAPSS results were known at that time. The protocol for N-CMAPSS, including four hypotheses, was written and hashed (SHA-256 7d75c109a05c0551\ldots{}, 2026-09-30T10:52:35Z) before the model or any competitor was run on those data (Section 5.4). N-CMAPSS had been used earlier in this project for a different question: fixed-time predictions of gradient-boosted, recurrent, convolutional, hidden-Markov and age-only models were known when the thresholds of the hypotheses were set, but no warning rule had ever been fitted or evaluated on it. The battery protocol (SHA-256 2a93fc12836f85d9\ldots{}) and the bearing protocol (6120e68d8aa34326\ldots{}) were written after the previous confirmation results were known and copy the model, the grids, the window and the four criteria unchanged (Section 5.4).

All evaluations split by engine. The main results use five-fold outer cross-validation (seed 2026); every engine is a test engine exactly once. Settings are selected with the inner cross-validation of Section 4.5.

\subsection{Competitors under an identical protocol}

The PHM08 challenge leaderboard used withheld test answers and a different task, so it cannot be compared directly. Instead we re-implemented the three winning challenge approaches and a convolutional baseline and ran them through exactly the same first-warning rule:

\begin{itemize}[leftmargin=1.2em,itemsep=2pt]
\item similarity-based prediction in the spirit of Wang et al. \cite{wang2008}: the median of a remaining-life distribution obtained by matching the recent principal-component trajectory against all training trajectories;
\item a recurrent network in the spirit of Heimes \cite{heimes2008}: a two-layer LSTM \cite{zheng2017} on 30-cycle windows of the seven temperature and shaft-speed sensors (s2, s3, s4, s8, s9, s13, s14), target capped at 125 cycles;
\item a multilayer perceptron with a Kalman filter over time in the spirit of Peel \cite{peel2008}: a single network (not an ensemble) with 64 and 32 hidden units on our causal features of the 15 sensors, target capped at 125, followed by a scalar Kalman filter in which the estimate decreases by one cycle per cycle, with the process-to-measurement noise ratio chosen with the other settings;
\item a convolutional network following Li et al. \cite{li2018}, on the same 30-cycle windows of seven sensors as the LSTM.
\end{itemize}

In this comparison every method, including ours, was fitted on the 75\% fit engines of each outer training part (seed 7) and had its warning settings (\textit{m}, \textit{k}, \textit{s} and its own option) chosen on the remaining 25\% calibration engines, with a shift grid of \textminus{}15 to +10. Our model therefore scores slightly lower here than in the main nested protocol, which uses all training engines. The methods are re-implementations in the spirit of the published approaches, not the original code, and they differ in their inputs: the LSTM and CNN see seven sensors, the MLP sees our encoder of 15 sensors.

To separate the architecture from the inputs and the target, three deep networks were also run with exactly what our model receives: the 15 regime-normalised sensors and the age, and a target cap chosen from \{40, 60, 90\} together with the other warning settings. They are the LSTM and the CNN above on 30-cycle windows of these inputs, and a hybrid Transformer with two self-attention blocks over the last 40 cycles of the sensors and regime indicators, concatenated with our causal features and trained with a quantile (pinball) loss, of which the median is used. All other steps of the protocol are unchanged.

Two further baselines were added after the main results were known, to test what the in-window rate measures. (i) A \textbf{constant-estimate control} keeps the warning times of our model but replaces its reported estimate by one constant per split, chosen on the training engines' out-of-fold warnings to maximise the number in the window. If the rate depended on the accuracy of the estimate, this control should be clearly worse. (ii) A \textbf{health-indicator threshold alarm} has no remaining-life regressor: a linear health indicator, in the spirit of Wang et al. \cite{wang2008}, is fitted by least squares of the exponentially weighted (span 10) baseline-removed sensors of Section 4.2 onto 0 for the first 30 cycles and 1 for the last five cycles of each training engine; the alarm is raised when the smoothed indicator stays at or above a threshold $\theta$ for \textit{k} cycles, and the reported estimate is a constant \textit{c} with \textit{H} \textminus{} 6 $\le$ \textit{c} $\le$ \textit{H}, the range of estimates our model reports at level \textit{H}. Its smoothing span, debounce, threshold (120 quantiles of the training indicator) and \textit{c} are chosen by the same nested out-of-fold procedure. Both baselines use the same folds, seeds and inner splits as our model, whose results the same runs reproduce engine by engine on all seven datasets.

\subsection{Three pre-specified confirmations}

Three datasets were analysed only after a protocol had been written and hashed (SHA-256 of the protocol text): N-CMAPSS (2026-09-30T10:52:35Z), the battery data (2026-09-30T12:54:24Z, before the data had been obtained for this project) and rolling bearings (2026-10-01T08:17:36Z, before any bearing measurement had been read; the protocol was also shared outside the analysis environment, in a time-stamped message, before the download). Each protocol fixed the data, any time scale, the model and its grids (unchanged), two competitors under the same tuning (an age-only regressor with the cycle count as its only input, and the MLP with Kalman filter of Section 5.3), the splits (primary: five-fold cross-validation by unit, seed 2026), the warning levels (primary \textit{H} = 20) and the same four hypotheses at \textit{H} = 20: in-window rate $\ge$ 85\%; late-or-missed rate $\le$ 10\%; higher in-window rate than age-only; and higher in-window rate than the MLP with Kalman filter (both by the exact McNemar test, two-sided \textit{p} < 0.05). The protocols were written in sequence, each after the previous result was known, but the model was not changed in response to any of them. The hashes were kept by the authors and not deposited with a registry, so the times are attested by the authors and by file time stamps only; a hash shows that a protocol has not changed since it was hashed, not when it was hashed. Recorded deviations are listed in the protocol files and in Supplementary Sections S5--S7.

\textbf{N-CMAPSS} (99 simulated engines) uses six temperature and shaft-speed channels and the three flight classes as regimes. \textbf{Battery} (123 commercial LFP/graphite cells; \cite{severson2019}) uses the seven per-cycle quantities of Section 4.1 as one regime. Battery lives (median 731 cycles) are about 3.5 times longer than PHM08 engine lives, so the protocol rescales time instead of re-tuning: consecutive cycles are averaged in blocks of $\kappa$ = max(1, round(median life / 209.5)) = 3, and the model, warning level and window are applied in blocks (in cycles the window is [\textminus{}39, +6] and \textit{H} = 20 is 60 cycles). Ten split-test records and five cells that never reach end of life were excluded before any model was run (Supplementary Section S5). \textbf{Bearings} are the 17 ball bearings run to failure on the PRONOSTIA platform for the IEEE PHM 2012 challenge \cite{nectoux2012} under three load and speed conditions, which serve as regimes. Each 0.1-s vibration snapshot, recorded every 10 s, is reduced to eight features (root mean square, peak, kurtosis and crest factor of the horizontal and vertical acceleration), and the same time-scale rule gives blocks of $\kappa$ = 7 snapshots (70 s). The bearing protocol also fixed, before it was run anywhere, a theory-guided shift rule (STW-T) that is described in Section 5.6.

\subsection{Leakage checks and statistics}

Four checks were run on PHM08: disjoint outer and inner engine sets (E1); causality, by recomputing predictions from truncated histories (E2); a negative control trained on labels built from permuted lives (E3); and exact reproduction of a refitted fold (E4) (Supplementary Section S4). The decision rule never uses the engine's total life, and the evaluation point is chosen by the model itself, so the leak inherent in evaluating at a cut point derived from the total life cannot arise. Proportions carry Clopper--Pearson 95\% intervals \cite{clopper1934}. Paired comparisons use the exact McNemar test \cite{mcnemar1947}, with Holm's correction over the seven competitors at each level of Table 4 and over the nine ablation variants of Table 5 \cite{holm1979}; the hypothesis tests of the confirmations are unadjusted, as pre-specified.

\subsection{Secondary analyses}

For the fixed-time reference analysis of Section 6.1, eleven models were evaluated at fixed fractions of each test engine's life in the same outer folds (Supplementary Section S1); the single models are scored on their raw estimates, without a window-centring shift, so their rates understate what a shift could add, whereas the stacked hybrid includes a tuned shift. Following the evaluation practice of the group that set the tolerance window, we also drew 50 random test engines and 168 training engines thirty times on PHM08 (seeds 1001--1030; inner-split seeds 2001--2030), with the same nested selection inside the 168 engines.

Three further analyses were added after the main results were known. (i) \textbf{Other windows} [\textminus{}10, +2], [\textminus{}20, +5] and [\textminus{}30, +10], with the settings re-selected for each window by the same nested procedure (Supplementary Section S8). (ii) \textbf{Maintenance cost}: in a renewal--reward model, an alarm triggers a planned replacement $\ell$ $\in$ \{0, 5, 10, 15\} cycles later; a unit that is never warned, or whose true remaining life at the alarm is at most $\ell$, fails. A planned replacement costs 1 and a failure $\rho$ $\in$ \{5, 10, 20, 50\}; the long-run cost per cycle is the total cost divided by the total cycles of use \cite{barlow1960}. Four alarm policies are compared, each tuned on the training units' out-of-fold predictions for every ($\ell$, $\rho$): our model tuned to the window, with the warning level \textit{H} $\in$ \{10, 15, 20, 30, 40\} chosen by cost (\textit{window-tuned}); our model tuned directly to the cost, with cap, span, debounce and an alarm threshold $\tau$ $\in$ \{1, \ldots{}, 60\} on the smoothed estimate chosen to minimise the training cost rate (\textit{cost-tuned}); the health-indicator alarm with span, debounce and threshold chosen the same way; and optimal age replacement with the age chosen on training lives. Run-to-failure and perfect foresight (replacement one cycle before failure) bound the range. On the batteries $\ell$ is in blocks (0--45 cycles) and \textit{H} is chosen like on the engines, unlike in the pre-specified battery protocol. Paired 95\% intervals for ratios of cost rates come from 2,000 bootstrap resamples of the test units, with the fitted policies fixed. (iii) \textbf{Theory-guided shift (STW-T)}: following Proposition 1, the shift is chosen so that the median out-of-fold warning error equals the window centre, \textminus{}5.5; it was pre-specified in the bearing protocol and is descriptive on the other datasets (Supplementary Section S10). The thirty random 50-engine test sets are reported in Supplementary Section S11.

\section{Results}\label{sec:6}

\subsection{Fixed-time prediction cannot meet the window}

\begin{figure*}[!tbp]\centering
\includegraphics[width=\textwidth]{img/Fig2_fixed_time_ceiling.png}
\caption{Fixed-time prediction on PHM08 (218 engines, five-fold cross-validation), scored in the window [\textminus{}13, +2]. (a) Share of predictions in the window when the prediction is made at fixed fractions of life (raw estimates). (b) The eleven reference models at half of life. In both panels the dashed line marks the self-timed warning model at \textit{H} = 20 (91.7\%).}
\end{figure*}

Fig. 2 shows that fixed-time prediction far from failure cannot meet the requirement. At half of life the best single model was the LSTM, with 30.3\% of predictions in the window. The physics-informed Bayesian model had the lowest mean absolute error (15.8 cycles) but 26.6\% in the window, and the stacked hybrid reached 33.5\%. Trained on one information source at a time, gradient-boosted trees placed 10.6\% in the window with age alone and 21.6\% with all sensor features. The share rises towards the end of life, to 64.7\% at 90\% of life for gradient-boosted trees, because wear then becomes visible; a fixed fraction of life is, however, not known in service. This motivates a policy that waits for the evidence.

\subsection{Main results}

\begin{table*}[!tbp]
\footnotesize\setlength{\tabcolsep}{4pt}
\caption{First-warning results of the self-timed warning model on the development and additional evaluation datasets (five-fold cross-validation, nested selection of the settings). In window: \textminus{}13 $\le$ \textit{e} $\le$ +2.}
\begin{tabular}{@{}>{\raggedright\arraybackslash}p{\dimexpr0.1274\linewidth-2\tabcolsep\relax}>{\raggedright\arraybackslash}p{\dimexpr0.0554\linewidth-2\tabcolsep\relax}>{\raggedright\arraybackslash}p{\dimexpr0.1053\linewidth-2\tabcolsep\relax}>{\raggedright\arraybackslash}p{\dimexpr0.0776\linewidth-2\tabcolsep\relax}>{\raggedright\arraybackslash}p{\dimexpr0.1274\linewidth-2\tabcolsep\relax}>{\raggedright\arraybackslash}p{\dimexpr0.0997\linewidth-2\tabcolsep\relax}>{\raggedright\arraybackslash}p{\dimexpr0.0886\linewidth-2\tabcolsep\relax}>{\raggedright\arraybackslash}p{\dimexpr0.1746\linewidth-2\tabcolsep\relax}>{\raggedright\arraybackslash}p{\dimexpr0.1440\linewidth-2\tabcolsep\relax}@{}}
\toprule
Dataset & \textit{H} & In window & \% & 95\% CI & Late or missed & Too early & Cycles left at warning, median (5--95\%) & Early limit can bind (\%) \\
\midrule
PHM08 & 10 & 210/218 & 96.3 & 92.9--98.4 & 4 & 4 & 10 (5--17) & 21.7 \\
 & 15 & 205/218 & 94.0 & 90.0--96.8 & 6 & 7 & 15 (8--22) & 65.6 \\
 & 20 & 200/218 & 91.7 & 87.3--95.0 & 8 & 10 & 19 (13--28) & 95.4 \\
 & 30 & 182/218 & 83.5 & 77.9--88.2 & 13 & 23 & 30 (20--43) & 100.0 \\
 & 40 & 173/218 & 79.4 & 73.4--84.5 & 22 & 23 & 38.5 (28--51) & 100.0 \\
FD001 & 10 & 98/100 & 98.0 & 93.0--99.8 & 2 & 0 & 11 (7--17) & 32.0 \\
 & 15 & 95/100 & 95.0 & 88.7--98.4 & 4 & 1 & 14 (8--21) & 59.0 \\
 & 20 & 88/100 & 88.0 & 80.0--93.6 & 4 & 8 & 19 (13--29) & 96.0 \\
FD002 & 10 & 250/260 & 96.2 & 93.0--98.1 & 5 & 5 & 11 (5--17) & 32.7 \\
 & 15 & 252/260 & 96.9 & 94.0--98.7 & 4 & 4 & 15 (9--22) & 73.1 \\
 & 20 & 246/260 & 94.6 & 91.1--97.0 & 7 & 7 & 19 (12--26) & 93.8 \\
FD003 & 10 & 94/100 & 94.0 & 87.4--97.8 & 5 & 1 & 11 (5--17) & 32.7 \\
 & 15 & 94/100 & 94.0 & 87.4--97.8 & 4 & 2 & 16 (8--24) & 77.8 \\
 & 20 & 96/100 & 96.0 & 90.1--98.9 & 1 & 3 & 21 (16--29) & 98.0 \\
FD004 & 10 & 225/249 & 90.4 & 86.0--93.7 & 7 & 17 & 10 (4--20) & 33.3 \\
 & 15 & 214/249 & 85.9 & 81.0--90.0 & 18 & 17 & 14 (6--26) & 61.8 \\
 & 20 & 206/249 & 82.7 & 77.5--87.2 & 17 & 26 & 20 (12--34) & 94.4 \\
\bottomrule
\end{tabular}\par\vspace{3pt}{\raggedright\scriptsize Pooled over PHM08 and FD001--FD004: 877/927 (94.6\%) at \textit{H} = 10, 860/927 (92.8\%) at \textit{H} = 15, 836/927 (90.2\%) at \textit{H} = 20. ``Early limit can bind'': share of warnings issued with at least 13 true cycles left. The confirmations are in Table 7.\par}
\end{table*}

Table 2 gives the main results. At the primary level \textit{H} = 20, 200 of 218 PHM08 engines (91.7\%, 95\% CI 87.3--95.0\%) received a correct first warning; 8 were late, 10 too early and none were missed. The warning came a median of 19 cycles before failure (5--95\%: 13--28). On C-MAPSS the rate was 88.0\% on FD001, 94.6\% on FD002, 96.0\% on FD003 and 82.7\% on FD004, the dataset with two fault modes and six regimes; pooled over the 927 engines it was 90.2\%. Lowering the warning level raised the rate (91.7\% $\rightarrow$ 96.3\% at \textit{H} = 10 on PHM08) at the cost of notice, and raising it to 30 and 40 lowered it to 83.5\% and 79.4\% (Fig. 3). Example trajectories are shown in Supplementary Fig. S1. All four leakage checks passed (Supplementary Section S4).

\begin{figure}[!tbp]\centering
\includegraphics[width=\linewidth]{img/Fig4_level_tradeoff.png}
\caption{Share of first warnings in the window as a function of the warning level \textit{H}. Error bars: 95\% Clopper--Pearson intervals for the pre-specified N-CMAPSS confirmation.}
\end{figure}

\subsection{What the in-window rate measures}

A policy could meet the window trivially by warning only in the last few cycles, where the early limit cannot bind. At \textit{H} = 10 this is largely what happens on PHM08: only 21.7\% of warnings came with at least 13 cycles left. At \textit{H} = 20 the share was 95.4\% on PHM08 and at least 93.8\% on every C-MAPSS dataset, and among those warnings the in-window rate was 92.8\% on PHM08 (Fig. 4b). The errors were concentrated around the centre of the window: the median error on PHM08 was \textminus{}4.5 cycles (interquartile range \textminus{}6.7 to \textminus{}1.7; Fig. 4a), close to the window centre of \textminus{}5.5, as the shift intends. Gross false alarms, more than 30 cycles early, did not occur at \textit{H} = 20 on PHM08, FD001 or FD003; there was one on FD002, and 6 of the 26 too-early warnings on FD004 were gross.

\begin{figure*}[!tbp]\centering
\includegraphics[width=\textwidth]{img/Fig5_error_leadtime.png}
\caption{Distribution of (a) the warning error and (b) the lead time (true cycles left) at \textit{H} = 20 for PHM08 (five-fold cross-validation) and N-CMAPSS (pre-specified confirmation). Shaded: the window; dotted: 13 cycles, below which the early limit practically cannot bind.}
\end{figure*}

\begin{table*}[!tbp]
\footnotesize\setlength{\tabcolsep}{4pt}
\caption{What the in-window rate measures (\textit{H} = 20, five-fold cross-validation). Our model against the same warning times with a constant reported estimate, and against a health-indicator threshold alarm without a remaining-life regressor.}
\begin{tabular}{@{}>{\raggedright\arraybackslash}p{\dimexpr0.1662\linewidth-2\tabcolsep\relax}>{\raggedright\arraybackslash}p{\dimexpr0.1385\linewidth-2\tabcolsep\relax}>{\raggedright\arraybackslash}p{\dimexpr0.1385\linewidth-2\tabcolsep\relax}>{\raggedright\arraybackslash}p{\dimexpr0.1385\linewidth-2\tabcolsep\relax}>{\raggedright\arraybackslash}p{\dimexpr0.1440\linewidth-2\tabcolsep\relax}>{\raggedright\arraybackslash}p{\dimexpr0.1440\linewidth-2\tabcolsep\relax}>{\raggedright\arraybackslash}p{\dimexpr0.1303\linewidth-2\tabcolsep\relax}@{}}
\toprule
Dataset & Our model: in window (\%) & Constant estimate: in window (\%) & Only ours / only constant; \textit{p} & Health-indicator alarm: in window (\%) & Only ours / only health indicator; \textit{p} & Health-indicator alarm: early limit can bind (\%) \\
\midrule
PHM08 & 200/218 (91.7) & 200 (91.7) & 7 / 7; 1.000 & 122 (56.0) & 83 / 5; < 0.001 & 81.7 \\
FD001 & 88/100 (88.0) & 90 (90.0) & 2 / 4; 0.688 & 51 (51.0) & 41 / 4; < 0.001 & 90.0 \\
FD002 & 246/260 (94.6) & 243 (93.5) & 9 / 6; 0.607 & 142 (54.6) & 105 / 1; < 0.001 & 80.8 \\
FD003 & 96/100 (96.0) & 94 (94.0) & 4 / 2; 0.688 & 57 (57.0) & 40 / 1; < 0.001 & 80.8 \\
FD004 & 206/249 (82.7) & 215 (86.3) & 5 / 14; 0.064 & 123 (49.4) & 94 / 11; < 0.001 & 78.2 \\
N-CMAPSS & 68/99 (68.7) & 70 (70.7) & 7 / 9; 0.804 & 53 (53.5) & 26 / 11; 0.020 & 77.8 \\
Battery & 120/123 (97.6) & 117 (95.1) & 4 / 1; 0.375 & 64 (52.0) & 57 / 1; < 0.001 & 85.4 \\
Engines pooled (PHM08, FD001--FD004) & 836/927 (90.2) & 842 (90.8) & 27 / 33; 0.519 & 495 (53.4) & -- & -- \\
\bottomrule
\end{tabular}\par\vspace{3pt}{\raggedright\scriptsize Exact McNemar tests (discordant units: only ours correct / only the control correct). Battery in blocks of three cycles. All levels \textit{H} = 10--40 are in Supplementary Section S12.\par}
\end{table*}

A warning is correct exactly when the true remaining life at the alarm lies between $\hat{y}$ \textminus{} 2 and $\hat{y}$ + 13 cycles. Because the alarm is raised when the smoothed, shifted estimate first stays below \textit{H}, the reported estimate varies little between units (standard deviation about 2 cycles on PHM08 at \textit{H} = 20, against about 5 cycles for the true remaining life at the alarm), so the rate may depend mainly on \textit{when} the alarm is raised. Table 3 and Supplementary Fig. S3 test this. With the warning times of our model and one constant estimate per split, the in-window rate at \textit{H} = 20 was 842/927 (90.8\%) over the five engine datasets, against 836/927 (90.2\%) with the model's own estimate (27 against 33 discordant engines, \textit{p} = 0.519); on PHM08 the two were identical (200 of 218). Over all 35 combinations of dataset and level from \textit{H} = 10 to 40, the constant estimate was never significantly worse; in one, N-CMAPSS at the exploratory level \textit{H} = 40, it was significantly better (21 against 7 discordant engines, unadjusted \textit{p} = 0.013). The in-window rate is therefore a measure of alarm timing: of how often the alarm falls in a band roughly 13 to 28 cycles before failure when \textit{H} = 20, not of how accurate the number attached to it is.

This does not make the regressor dispensable. The health-indicator threshold alarm, which times its alarm on a classical linear health indicator instead of on the estimated remaining life, placed only 49.4--57.0\% of first warnings in the window at \textit{H} = 20 (495/927 = 53.4\% over the engines), significantly below our model on every dataset, with many alarms too early (64 of 218 on PHM08). What the horizon-focused regressor contributes is a health indicator calibrated in cycles to failure, which lets a single threshold fire at a similar remaining life on units that degrade at different rates. Proposition 1 should be read the same way: the concentration that limits the rate is the concentration of the alarm time relative to failure.

\subsection{Comparison under an identical protocol}

\begin{table*}[!tbp]
\footnotesize\setlength{\tabcolsep}{4pt}
\caption{Identical first-warning protocol on PHM08 (218 engines; fit on 75\% of each outer training part, warning settings chosen on the other 25\%). Upper rows of each level: re-implemented challenge approaches and CNN with their original inputs; last three rows: deep networks given exactly the inputs and target options of our model.}
\begin{tabular}{@{}>{\raggedright\arraybackslash}p{\dimexpr0.2493\linewidth-2\tabcolsep\relax}>{\raggedright\arraybackslash}p{\dimexpr0.0499\linewidth-2\tabcolsep\relax}>{\raggedright\arraybackslash}p{\dimexpr0.1274\linewidth-2\tabcolsep\relax}>{\raggedright\arraybackslash}p{\dimexpr0.1163\linewidth-2\tabcolsep\relax}>{\raggedright\arraybackslash}p{\dimexpr0.0942\linewidth-2\tabcolsep\relax}>{\raggedright\arraybackslash}p{\dimexpr0.0886\linewidth-2\tabcolsep\relax}>{\raggedright\arraybackslash}p{\dimexpr0.1551\linewidth-2\tabcolsep\relax}>{\raggedright\arraybackslash}p{\dimexpr0.1192\linewidth-2\tabcolsep\relax}@{}}
\toprule
Method & \textit{H} & In window (\%) & 95\% CI & Late or missed & Too early & Only ours correct / only competitor correct & Holm-adjusted \textit{p} \\
\midrule
Self-timed warning (ours) & 10 & 206 (94.5) & 90.6--97.1 & 9 & 3 & -- & -- \\
MLP + Kalman filter (Peel-style) & 10 & 195 (89.4) & 84.6--93.2 & 9 & 14 & 22 / 11 & 0.240 \\
LSTM (Heimes-style) & 10 & 173 (79.4) & 73.4--84.5 & 16 & 29 & 42 / 9 & < 0.001 \\
Similarity (Wang-style) & 10 & 160 (73.4) & 67.0--79.1 & 21 & 37 & 55 / 9 & < 0.001 \\
CNN \cite{li2018} & 10 & 138 (63.3) & 56.5--69.7 & 40 & 40 & 78 / 10 & < 0.001 \\
LSTM, same inputs and cap set as ours & 10 & 201 (92.2) & 87.8--95.4 & 10 & 7 & 11 / 6 & 0.665 \\
CNN, same inputs and cap set as ours & 10 & 179 (82.1) & 76.4--87.0 & 20 & 19 & 35 / 8 & < 0.001 \\
Hybrid Transformer, same inputs and cap set as ours & 10 & 206 (94.5) & 90.6--97.1 & 10 & 2 & 8 / 8 & 1.000 \\
Self-timed warning (ours) & 20 & 195 (89.4) & 84.6--93.2 & 10 & 13 & -- & -- \\
MLP + Kalman filter (Peel-style) & 20 & 180 (82.6) & 76.9--87.4 & 15 & 23 & 30 / 15 & 0.071 \\
LSTM (Heimes-style) & 20 & 140 (64.2) & 57.5--70.6 & 29 & 49 & 70 / 15 & < 0.001 \\
Similarity (Wang-style) & 20 & 138 (63.3) & 56.5--69.7 & 34 & 46 & 64 / 7 & < 0.001 \\
CNN \cite{li2018} & 20 & 121 (55.5) & 48.6--62.2 & 46 & 51 & 82 / 8 & < 0.001 \\
LSTM, same inputs and cap set as ours & 20 & 176 (80.7) & 74.9--85.7 & 27 & 15 & 36 / 17 & 0.038 \\
CNN, same inputs and cap set as ours & 20 & 172 (78.9) & 72.9--84.1 & 22 & 24 & 40 / 17 & 0.013 \\
Hybrid Transformer, same inputs and cap set as ours & 20 & 201 (92.2) & 87.8--95.4 & 5 & 12 & 13 / 19 & 0.377 \\
\bottomrule
\end{tabular}\par\vspace{3pt}{\raggedright\scriptsize Exact McNemar tests against our model; Holm correction over the seven competitors at each level.\par}
\end{table*}

Table 4 and Fig. 5a compare the methods. Against the four re-implementations with their original inputs, our model had the highest in-window rate at both levels: 89.4\% at \textit{H} = 20, against 82.6\% for the MLP with Kalman filter, 64.2\% for the LSTM, 63.3\% for the similarity method and 55.5\% for the CNN. The differences from the LSTM, similarity and CNN were significant at both levels after Holm's correction (\textit{p} < 0.001). Against the MLP with Kalman filter, which shares our feature encoder, the discordant engines were 30 against 15 at \textit{H} = 20 (unadjusted \textit{p} = 0.036, Holm-adjusted 0.071) and 22 against 11 at \textit{H} = 10 (Holm-adjusted 0.240); neither difference is significant after correction.

The lower block of Table 4 removes the differences in inputs and target. Given the same 15 regime-normalised sensors and the age as our model, and choosing their target cap from the same set, the networks improved markedly: at \textit{H} = 20 the LSTM rose from 64.2\% to 80.7\% and the CNN from 55.5\% to 78.9\%. Both remained significantly below our model at \textit{H} = 20 (Holm-adjusted \textit{p} = 0.038 and 0.013); at \textit{H} = 10 only the CNN did. The hybrid Transformer, which also receives our causal features, reached 201/218 (92.2\%) at \textit{H} = 20, slightly more than our model's 195/218, and 206/218 (94.5\%) at \textit{H} = 10, the same as ours; neither difference was significant (13 against 19 discordant engines at \textit{H} = 20, \textit{p} = 0.377). Much of the gap to the original re-implementations is therefore due to their inputs and targets; with our inputs, one of the three networks, the Transformer, matched gradient boosting.

\begin{figure*}[!tbp]\centering
\includegraphics[width=\textwidth]{img/Fig6_comparison_ablation.png}
\caption{(a) In-window rate under the identical protocol of Table 4 with 95\% intervals. (b) Component ablation on PHM08 at \textit{H} = 20 (nested protocol); each bar changes one aspect of the model (the untuned rule fixes all four settings; age only replaces all inputs), and all remaining settings are re-tuned.}
\end{figure*}

\subsection{Component ablation}

\begin{table*}[!tbp]
\footnotesize\setlength{\tabcolsep}{4pt}
\caption{Component ablation on PHM08 (218 engines, nested protocol, \textit{H} = 20; \textit{H} = 10 in the last column).}
\begin{tabular}{@{}>{\raggedright\arraybackslash}p{\dimexpr0.2991\linewidth-2\tabcolsep\relax}>{\raggedright\arraybackslash}p{\dimexpr0.0942\linewidth-2\tabcolsep\relax}>{\raggedright\arraybackslash}p{\dimexpr0.0665\linewidth-2\tabcolsep\relax}>{\raggedright\arraybackslash}p{\dimexpr0.0886\linewidth-2\tabcolsep\relax}>{\raggedright\arraybackslash}p{\dimexpr0.0776\linewidth-2\tabcolsep\relax}>{\raggedright\arraybackslash}p{\dimexpr0.0997\linewidth-2\tabcolsep\relax}>{\raggedright\arraybackslash}p{\dimexpr0.1662\linewidth-2\tabcolsep\relax}>{\raggedright\arraybackslash}p{\dimexpr0.1081\linewidth-2\tabcolsep\relax}@{}}
\toprule
Variant & In window & \% & Late or missed & Too early & Cycles left, median & Only full / only variant correct; Holm \textit{p} & \% in window at \textit{H} = 10 \\
\midrule
Full model & 200/218 & 91.7 & 8 & 10 & 19 & -- & 96.3 \\
Conventional target cap 125 & 196/218 & 89.9 & 10 & 12 & 19 & 10 / 6; 1.000 & 96.8 \\
No target cap & 178/218 & 81.7 & 22 & 18 & 19 & 29 / 7; 0.002 & 92.7 \\
No window centring (\textit{s} = 0) & 168/218 & 77.1 & 49 & 1 & 19 & 41 / 9; < 0.001 & 78.0 \\
No debounce (\textit{k} = 1) & 195/218 & 89.4 & 8 & 15 & 23 & 13 / 8; 1.000 & 94.5 \\
No smoothing (\textit{m} = 1) & 197/218 & 90.4 & 12 & 9 & 19 & 10 / 7; 1.000 & 95.4 \\
Untuned rule (\textit{c} = 60, \textit{m} = 1, \textit{k} = 1, \textit{s} = 0) & 171/218 & 78.4 & 46 & 1 & 20 & 39 / 10; < 0.001 & 83.9 \\
Seven temperature and speed sensors (with age) & 183/218 & 83.9 & 22 & 13 & 19 & 25 / 8; 0.023 & 94.5 \\
Sensor features without age & 187/218 & 85.8 & 17 & 14 & 18 & 17 / 4; 0.029 & 96.3 \\
Age only & 31/218 & 14.2 & 130 & 57 & 31.5 & 169 / 0; < 0.001 & 6.0 \\
Oracle-timed control (not deployable) & 191/218 & 87.6 & 22 & 5 & 20 & 19 / 10; 0.136 & 95.9 \\
\bottomrule
\end{tabular}\par\vspace{3pt}{\raggedright\scriptsize Paired comparison with the full model by the exact McNemar test, Holm-adjusted over the nine variants (unadjusted for the oracle-timed control). Oracle-timed control: the full model's calibrated estimate evaluated at the cycle where the true remaining life equals \textit{H}; this uses the true life and cannot be deployed.\par}
\end{table*}

Table 5 and Fig. 5b isolate the components on PHM08 at \textit{H} = 20. The full model reproduced the main result exactly (200/218). Removing the window-centring shift caused the largest loss: 168 engines (77.1\%), with 49 late warnings, because about half of the unshifted errors are positive and every error above +2 is late, while the wide early side of the window goes unused (Holm-adjusted \textit{p} < 0.001). An untuned rule (cap 60, no smoothing, no debounce, no shift) reached 78.4\%. Training without a target cap gave 81.7\% (\textit{p} = 0.002), but the conventional cap of 125 gave 89.9\%, statistically indistinguishable from the selected caps (10 against 6 discordant engines, \textit{p} = 1.000). Smoothing (90.4\% without) and debouncing (89.4\% without) made no significant difference. The inputs mattered: the seven temperature and shaft-speed sensors (with age) gave 83.9\% (\textit{p} = 0.023), sensor features without the age gave 85.8\% (\textit{p} = 0.029), and age alone gave 14.2\%, with 104 engines never warned. Finally, the oracle-timed control, which evaluates the same calibrated estimate at the cycle where exactly 20 true cycles remain, reached 87.6\%. Letting the model choose its own moment therefore did not cost accuracy relative to being told the true moment (19 against 10 discordant engines, \textit{p} = 0.136).

\subsection{Maintenance cost of the alarms}

\begin{table*}[!tbp]
\footnotesize\setlength{\tabcolsep}{4pt}
\caption{Long-run cost rate relative to perfect foresight (1.00 = replacement one cycle before every failure) for alarm policies tuned on training units, at two failure-cost ratios $\rho$ and lead times $\ell$ of 0 and 10 cycles (blocks for the battery). STW (window): our model tuned to the window, warning level chosen by cost; STW (cost): our model tuned directly to the cost; HI (cost): health-indicator threshold alarm tuned to the cost; Age: optimal age replacement.}
\begin{tabular}{@{}>{\raggedright\arraybackslash}p{\dimexpr0.1322\linewidth-2\tabcolsep\relax}>{\raggedright\arraybackslash}p{\dimexpr0.0386\linewidth-2\tabcolsep\relax}>{\raggedright\arraybackslash}p{\dimexpr0.1047\linewidth-2\tabcolsep\relax}>{\raggedright\arraybackslash}p{\dimexpr0.1047\linewidth-2\tabcolsep\relax}>{\raggedright\arraybackslash}p{\dimexpr0.1047\linewidth-2\tabcolsep\relax}>{\raggedright\arraybackslash}p{\dimexpr0.1047\linewidth-2\tabcolsep\relax}>{\raggedright\arraybackslash}p{\dimexpr0.1047\linewidth-2\tabcolsep\relax}>{\raggedright\arraybackslash}p{\dimexpr0.1047\linewidth-2\tabcolsep\relax}>{\raggedright\arraybackslash}p{\dimexpr0.1047\linewidth-2\tabcolsep\relax}>{\raggedright\arraybackslash}p{\dimexpr0.0965\linewidth-2\tabcolsep\relax}@{}}
\toprule
Dataset & $\ell$ & $\rho$ = 10: STW (window) & STW (cost) & HI (cost) & Age & $\rho$ = 50: STW (window) & STW (cost) & HI (cost) & Age \\
\midrule
PHM08 & 0 & 1.09 & 1.12 & 1.14 & 1.77 & 1.28 & 1.50 & 1.33 & 2.37 \\
 & 10 & 1.10 & 1.17 & 1.12 & 1.77 & 1.10 & 1.75 & 1.12 & 2.37 \\
FD001 & 0 & 1.06 & 1.14 & 1.28 & 1.74 & 1.06 & 1.56 & 2.15 & 2.38 \\
 & 10 & 1.10 & 1.15 & 1.21 & 1.74 & 1.10 & 1.57 & 1.65 & 2.38 \\
FD002 & 0 & 1.06 & 1.08 & 1.15 & 1.67 & 1.06 & 1.24 & 1.32 & 1.92 \\
 & 10 & 1.09 & 1.16 & 1.19 & 1.67 & 1.09 & 1.64 & 1.37 & 1.92 \\
FD003 & 0 & 1.08 & 1.14 & 1.20 & 1.86 & 1.08 & 1.56 & 1.63 & 2.54 \\
 & 10 & 1.06 & 1.04 & 1.21 & 1.86 & 1.06 & 1.04 & 1.65 & 2.54 \\
FD004 & 0 & 1.11 & 1.13 & 1.10 & 2.13 & 1.28 & 1.47 & 1.10 & 2.30 \\
 & 10 & 1.16 & 1.22 & 1.20 & 2.13 & 1.36 & 1.91 & 1.56 & 2.30 \\
N-CMAPSS & 0 & 1.79 & 1.51 & 2.33 & 1.72 & 1.63 & 2.07 & 3.64 & 2.35 \\
 & 10 & 1.70 & 1.51 & 3.01 & 1.72 & 2.50 & 2.08 & 2.44 & 2.35 \\
Battery & 0 & 1.13 & 1.13 & 1.06 & 2.23 & 1.55 & 1.47 & 1.06 & 4.59 \\
 & 10 & 1.13 & 1.13 & 1.18 & 2.23 & 1.54 & 1.48 & 1.53 & 4.59 \\
\bottomrule
\end{tabular}\par\vspace{3pt}{\raggedright\scriptsize Five-fold cross-validation. All 16 combinations of $\ell$ $\in$ \{0, 5, 10, 15\} and $\rho$ $\in$ \{5, 10, 20, 50\}, with paired bootstrap intervals, are in Supplementary Section S12. Mean regret over the 112 dataset--$\ell$--$\rho$ combinations, relative to the cheapest of the three alarm policies in each: window-tuned STW 5.5\% (median 0.0\%, maximum 73.5\%), cost-tuned STW 10.9\% (median 2.4\%, maximum 94.9\%), health-indicator alarm 22.2\%, age replacement 81.7\%.\par}
\end{table*}

Table 6 and Fig. 6 price the alarms. The window-tuned model, with its warning level chosen by cost on the training units, brought the long-run cost rate on PHM08 to 1.07--1.28 times that of perfect foresight without lead time, against 1.70--2.37 for optimal age replacement, and it was significantly cheaper than age replacement in 94 of the 96 combinations of lead time and cost ratio on the five engine datasets and the batteries and never significantly more expensive. On N-CMAPSS, where the alarm time was least concentrated (Section 6.8), it was not significantly cheaper than age replacement in any combination. Against the health-indicator alarm tuned directly for cost, it was significantly cheaper in 60 of 112 combinations and significantly more expensive in 4, all of them on the batteries, where a conservative health-indicator threshold avoided every failure at no lead time (0 failures over the four cost ratios).

Tuning our model directly for the cost did not improve on tuning it for the window. The cost-tuned policy was significantly cheaper in 33 of 112 combinations, by 9.7\% on average, but its losses were larger and harder to predict: up to 94.9\% above the best policy (for example, 1.75 against 1.10 times perfect foresight on PHM08 with $\ell$ = 10 and $\rho$ = 50), against at most 45.6\% for the window-tuned policy outside N-CMAPSS. In that example the cost-tuned alarms came a median of 19 cycles before failure against 30 for the window-tuned ones, and 3 test engines failed against 0: minimising the empirical cost over a large grid selects alarms that leave little margin on the training units, and with $\rho$ = 50 each additional test failure outweighs many cycles of use. The window, with its narrow late side, builds a margin into the tuning, and only the warning level adapts to the cost. Over all 112 combinations the window-tuned policy had the lowest mean regret (5.5\%, median 0.0\%), followed by the cost-tuned policy (10.9\%), the health-indicator alarm (22.2\%) and age replacement (81.7\%). In this sense the tolerance window is a robust surrogate for the maintenance cost rather than an arbitrary acceptance rule.

\begin{figure*}[!tbp]\centering
\includegraphics[width=\textwidth]{img/FigR2_cost_primary.png}
\caption{Long-run cost rate relative to perfect foresight as a function of the failure-to-replacement cost ratio $\rho$, without lead time (top) and with a lead time of 10 cycles (bottom; 10 blocks of three cycles for the battery), for our model tuned to the window (dashed) or to the cost (solid blue), a cost-tuned health-indicator alarm (orange) and optimal age replacement (dotted).}
\end{figure*}

\subsection{The three pre-specified confirmations}

\begin{table*}[!tbp]
\footnotesize\setlength{\tabcolsep}{4pt}
\caption{Pre-specified confirmations (five-fold cross-validation by unit, \textit{H} = 20) and hypothesis outcomes.}
\begin{tabular}{@{}>{\raggedright\arraybackslash}p{\dimexpr0.2105\linewidth-2\tabcolsep\relax}>{\raggedright\arraybackslash}p{\dimexpr0.2659\linewidth-2\tabcolsep\relax}>{\raggedright\arraybackslash}p{\dimexpr0.2659\linewidth-2\tabcolsep\relax}>{\raggedright\arraybackslash}p{\dimexpr0.2577\linewidth-2\tabcolsep\relax}@{}}
\toprule
Hypothesis (\textit{H} = 20) & N-CMAPSS (99 simulated engines) & Battery (123 cells, blocks of 3 cycles) & Bearings (17 bearings, blocks of 7 snapshots) \\
\midrule
1. In-window rate $\ge$ 85\% & 68/99 = 68.7\% (58.6--77.6\%); not supported & 120/123 = 97.6\% (93.0--99.5\%); supported & 0/17 = 0.0\% (0.0--19.5\%); not supported \\
2. Late or missed $\le$ 10\% & 18 + 5 = 23.2\%; not supported & 2 + 1 = 2.4\%; supported & 0 + 10 = 58.8\%; not supported \\
3. Better than age-only & 24.2\%; 49 / 5, \textit{p} < 0.001; supported & 0.8\%; 119 / 0, \textit{p} < 0.001; supported & 5.9\%; 0 / 1, \textit{p} = 1.000; not supported \\
4. Better than MLP + Kalman & 58.6\%; 25 / 15, \textit{p} = 0.154; not supported & 85.4\%; 17 / 2, \textit{p} < 0.001; supported & 0.0\%; 0 / 0, \textit{p} = 1.000; not supported \\
\bottomrule
\end{tabular}\par\vspace{3pt}{\raggedright\scriptsize Rows 1--2: count and rate of the self-timed model (95\% CI). Rows 3--4: competitor rate; discordant units (only ours correct / only competitor correct) and exact McNemar \textit{p}. Secondary analyses are in Supplementary Sections S5--S7.\par}
\end{table*}

Table 7 summarises the three confirmations, which span the range from success to complete failure. \textbf{Battery.} On measured battery degradation the model placed 120 of 123 first warnings in the window (97.6\%), with 2 late warnings, 1 missed cell (the shortest-lived, 148 cycles) and none too early; all four hypotheses were supported. The warnings came a median of 57 cycles before failure, and the early limit could bind for 96.7\% of them. Age-only never warned for 120 of 123 cells, because cell lives range from 148 to 1935 cycles. Fig. 7a shows the capacity fade: slow at first, then a sharp knee, on which the warnings fall. Trained on two test batches and tested on the third, a harder secondary test, the model reached 69.9\% (27 late), still above the MLP with Kalman filter (39.0\%; Supplementary Section S5).

\textbf{N-CMAPSS.} The model placed 68 of 99 first warnings in the window (68.7\%); 18 warnings were late and 5 engines were missed. It clearly outperformed age-only (hypothesis 3), but neither threshold was met and its advantage over the MLP with Kalman filter was not significant. An exploratory re-run with all 14 channels gave 65.7\%, no better (Supplementary Section S6).

\textbf{Bearings.} On the rolling bearings the model failed completely: none of the 17 first warnings fell in the window; 7 warnings were too early and 10 bearings were never warned. The competitors did no better (age-only 1 of 17, MLP with Kalman filter 0), so no hypothesis was supported. The pre-specified official challenge split (6 training, 11 test bearings) gave the same picture (Supplementary Section S7). A diagnosis made after the run, with an ad hoc criterion, explains part of this: in 11 of the 17 bearings the vibration level rose to and stayed above 1.5 times its early-life level only in the last six blocks (about seven minutes) before failure, whereas a correct warning at \textit{H} = 20 needs a signal about 20 blocks before failure (Supplementary Table S1). It does not explain everything: the model also missed or warned too early on bearings with long degradation, and raised early false alarms on several abruptly failing ones. Small numbers (13--14 training bearings per fold) and three different operating conditions are equally plausible causes. The pre-specified theory-guided shift (STW-T) was equally unsuccessful (0.0\%; 0 against 0 discordant bearings, \textit{p} = 1.000), so its non-inferiority criterion was met only trivially and the test was uninformative.

\begin{figure*}[!tbp]\centering
\includegraphics[width=\textwidth]{img/Fig7_battery.png}
\caption{Battery confirmation. (a) Discharge capacity of the 123 cells (grey) and the first warning of the self-timed model at \textit{H} = 20 blocks (primary cross-validation), coloured by outcome; the missed cell has no marker. (b) In-window rate of the three pre-specified methods in the primary analysis and two secondary analyses.}
\end{figure*}

\subsection{What limits the in-window rate}

Fig. 8 checks Proposition 1 on every combination of dataset and warning level. For each, we took the test warnings, recovered the unshifted error $\varepsilon$ = \textit{e} \textminus{} \textit{s}, and computed its median \textit{m} and its centred concentration Pr(|$\varepsilon$ \textminus{} \textit{m}| $\le$ 7.5); missed units count as failures. On the seven datasets other than the bearings, the realised in-window rate stayed close to the concentration (mean absolute difference 1.4 percentage points, largest 5.1), sometimes above it, because the empirical errors are not exactly symmetric and assumption A1 holds only approximately; the concentration is therefore a reference value rather than a strict bound, and since both quantities come from the same warnings their agreement is partly built in. The informative contrast is between datasets: at \textit{H} = 20 the concentration was 96.7\% on the batteries, 91.3\% on PHM08, 83.9\% on FD004, 72.7\% on N-CMAPSS, the same ordering as the realised rates; on the bearings, where 10 of 17 units were never warned, the concentration is not meaningful.

If the datasets differ in concentration rather than in how well the shift is set, a better shift rule should not help, and it did not: choosing the shift so that the median out-of-fold error equals the window centre (STW-T) changed the rate at \textit{H} = 20 by between \textminus{}2.0 and 0.8 percentage points on the seven datasets seen before it was written (Supplementary Section S10). Section 6.3 adds the practical reading: since the reported estimate hardly varies at the alarm, the concentration of $\varepsilon$ is essentially the concentration of the alarm time relative to failure. The rate also grew with the width of the window on every dataset (Supplementary Section S8).

\begin{figure*}[!tbp]\centering
\includegraphics[width=\textwidth]{img/Fig8_theory_check.png}
\caption{Check of Proposition 1 on the test warnings at \textit{H} = 10, 15 and 20. (a) Realised in-window rate against the centred concentration of the unshifted error, Pr(|$\varepsilon$ \textminus{} \textit{m}| $\le$ 7.5), the attainable rate under Proposition 1 for a symmetric unimodal error. (b) Shift selected on training data (mean over test units) against the optimal shift implied by the test errors (bearings omitted, since most were never warned). Dashed: identity. Battery and bearing values in blocks.}
\end{figure*}

\section{Discussion}\label{sec:7}

\subsection{Why self-timing works, and what it does not do}

The warning model is not a better long-range predictor; its regressor was designed for, and evaluated at, the end of life only. What it does is time an alarm (Section 6.3): the score rewards an alarm raised roughly 13 to 28 cycles before failure at \textit{H} = 20, and the regressor earns its place as a health indicator calibrated in cycles to failure, which a classical linear health indicator with a tuned threshold could not replace. The ablation agrees: removing the window-centring shift, which moves the alarm, cost the most, followed by removing the target cap and reducing the sensor set, whereas smoothing, debouncing and the exact cap value mattered little, and the oracle-timed control shows that letting the model choose its own moment did not cost accuracy. The same reading explains Section 6.6: the window tolerates 13 cycles of early error but only 2 of late error, so tuning for it builds a margin into the alarm, and choosing only the warning level by cost gave the lowest regret, whereas minimising the training cost directly was fragile when failures were expensive. The regressor itself is exchangeable (Section 6.4): with our inputs, a hybrid Transformer matched gradient boosting (92.2\% against 89.4\%) at about 36 minutes of training per fold against about 34 seconds, so gradient boosting remains our default for its speed and auditability. No claim about the PHM08 challenge leaderboard follows, because the task and data split differ. The model does not provide long-range planning information, for which fixed-time predictions and their uncertainty remain the relevant tools \cite{kamariotis2024,mitici2023}; a planning model sets the maintenance slot, and the warning model triggers the final removal.

\subsection{When the method works: lessons from the three confirmations}

The confirmations are the most informative results of this study, because their protocols were fixed before any warning result on those data was seen, and together they delimit the method. It works when the approach of failure becomes visible, consistently across units, some time before the warning level: battery capacity declines slowly and then drops at a knee, and the error at the warning was more concentrated than on any engine dataset (interquartile range 3.5 blocks). It degrades when that signal is noisy relative to the remaining time, as on N-CMAPSS, where short lives place a 20-cycle warning at about three quarters of life and only 72.7\% of the errors lay within $\pm$7.5 cycles of their median; neither more channels nor a theory-guided shift helped. It failed completely on the bearings, where most failures came within minutes of a sustained rise in vibration, although that post hoc diagnosis does not explain every bearing and the small number of units may matter as much. A shift between test campaigns also costs accuracy, as the battery leave-one-batch-out analysis showed (69.9\%). Proposition 1 gives a practical check before deployment: estimate the concentration of the out-of-fold error at the intended warning level on the fleet's own run-to-failure history; if it is low, no choice of shift will rescue the rate, and either the warning level, the window or the sensing must change.

\subsection{Limitations}

\begin{itemize}[leftmargin=1.2em,itemsep=2pt]
\item \textbf{No field fleet.} Six of the eight datasets are simulations, and the battery and bearing data are laboratory measurements. The confirmations range from 97.6\% to 0.0\%, so the attainable rate depends on the system and must be re-estimated for each application.
\item \textbf{Development on PHM08.} The model and the grids were fixed while looking at PHM08 results, and the primary level was designated after the C-MAPSS results were known, so these estimates may be optimistic; only the three confirmations were pre-specified, and their hashes were not deposited with a registry.
\item \textbf{Post hoc secondary analyses.} The constant-estimate control, the health-indicator alarm, the cost-tuned policies, the ablation, the oracle-timed control, the networks with equal inputs, the other windows, the cost analysis, the 14-channel N-CMAPSS run, the check of Proposition 1, STW-T on the seen datasets, the thirty random 50-engine tests and the bearing diagnosis were added after the main results were known. They describe the method; they are not confirmations.
\item \textbf{Competitors.} The challenge approaches are re-implementations in the spirit of the published methods, not the original code. The networks' architectures were taken from the literature and only their target cap was tuned; more extensive tuning could improve them.
\item \textbf{Cost-tuned comparison.} The cost-tuned policies choose among about two thousand settings, whereas the window-tuned policy chooses only its warning level by cost after the window tuning, so part of the fragility of direct cost tuning reflects selection over a larger grid on limited training units; a regularised or probabilistic cost tuning \cite{kamariotis2024} might close the gap. The health-indicator alarm is one classical baseline, not the strongest possible.
\item \textbf{A simple cost model.} The renewal--reward model assumes a fixed lead time, independent units and replacement as good as new, and the cost ratios are illustrative rather than measured; it ignores spare-part inventory, downtime and the grouping of maintenance across units \cite{depater2022,mitici2023}.
\end{itemize}

\section{Future work: test and maintenance policies, and automatic fault detection and diagnosis}\label{sec:8}

The warning model was developed and evaluated as a trigger for the replacement of single units. Its two ingredients, an alarm rule tuned on run-to-failure histories and a leakage-safe protocol for evaluating that rule, are not tied to this use. Two directions within the scope of this journal follow directly, and the present results indicate both why the model is suited to them and what remains to be shown (Fig. 9).

\begin{figure*}[!tbp]\centering
\includegraphics[width=\textwidth]{img/Fig9_future_work.png}
\caption{Planned uses of the self-timed warning model in test and maintenance policies (left) and in automatic fault detection and diagnosis (right). Solid boxes: evaluated in this paper; dashed boxes: future work.}
\end{figure*}

\subsection{Test and maintenance policies}

Condition-based replacement with a lead time is the first application, and Section 6.6 gives its long-run cost against age replacement and perfect foresight. Three extensions are natural. First, logistics: an alarm can order the spare part, so the warning level should be chosen jointly with the order lead time and the cost of a failure; Section 6.6 already chooses the level by cost for a given lead time, and a closed-form or dynamic-programming rule for this choice would make the policy transparent. Second, fleets: alarms of several units can be grouped into shared maintenance visits, as in alarm-based scheduling for aircraft engines \cite{depater2022} and dynamic maintenance of multiple components \cite{mitici2023}. How much slack such grouping has is set by the spread of the alarm time relative to failure, the quantity that Proposition 1 and Section 6.3 identify as the limit of the in-window rate. Third, imperfect repair, downtime and spare-part stocks, which our renewal--reward model leaves out.

The model can also set \textbf{test and inspection policies}, a long-standing topic of maintenance optimisation \cite{dejonge2020}, for units that are not monitored continuously, such as standby equipment that is proof-tested or structures that are inspected at intervals. There the alarm would trigger an additional test rather than a replacement, while a baseline test interval covers the remaining risk. Two pieces are missing: an encoder for sparse and irregular observations, since the present encoder assumes one observation per cycle, and a cost model with test costs, test-induced downtime and the chance that a test reveals a dormant failure. The asymmetric window has a direct reading in this setting, its early limit being the margin accepted between test and failure and its late limit the risk of testing too late, and Proposition 1 gives a check before such a policy is fielded: if the alarm time is not concentrated at the intended level, no threshold will make on-demand testing reliable, and the baseline interval must carry the risk. Because alarms tuned on one test campaign lost accuracy on another (battery leave-one-batch-out: 69.9\%), such policies should also include a fallback, for example age-based replacement or the baseline test interval for units outside the conditions seen in training.

\subsection{Automatic fault detection and diagnosis}

The trigger of Section 4.4 is a sequential decision rule, a debounced threshold crossing on a smoothed statistic, of the same kind as the control charts and run rules used for data-driven fault detection \cite{venkatasubramanian2003}. Section 6.3 showed that, under the asymmetric window, the score rewards the timing of the alarm, which is also what a fault detector trades off: detection delay against false alarms. A first study should therefore compare the warning model with established detectors, such as cumulative-sum \cite{page1954} and exponentially weighted moving-average charts \cite{roberts1959} and Hotelling's T\textsuperscript{2} statistic on the regime-normalised sensors, under both the window and the usual detection metrics (detection delay, false-alarm rate and missed detections), with the same nested, leakage-safe protocol. Gross false alarms, more than 30 cycles early, were rare in our results except on FD004 (6 of its 26 too-early warnings at \textit{H} = 20), which gives a reference for that comparison.

Diagnosis is the second step. FD003 and FD004 contain two fault modes, and FD004 had the lowest in-window rate of the C-MAPSS datasets. A diagnostic extension would estimate the fault mode together with the remaining life and issue mode-specific alarms, with the warning level and window set per mode. Because C-MAPSS does not label the fault mode of each engine, this requires data with labelled modes, such as N-CMAPSS, whose subsets simulate different failure modes. Abrupt faults are the third topic. On the bearings, the vibration level rose to and stayed above 1.5 times its early-life level only in the last six blocks before failure in 11 of 17 units, so no alarm timed from these features could give 20 blocks of notice. A two-stage design, in which a fast detector on high-frequency vibration features declares the onset of a defect and the warning model then times the replacement, is a plausible remedy that we have not tested. Finally, the regime-normalised, baseline-removed signals behave like residuals and could be used to flag sensor faults before they corrupt the alarm; the robustness of the warning model to sensor faults has not been examined.

For both directions, the evaluation practice of this paper carries over unchanged: nested selection by unit, controls that separate what a score measures from what it is meant to measure, and protocols hashed before confirmation data are analysed.

\section{Conclusion}\label{sec:9}

Requiring a remaining-life estimate to lie within \textminus{}13 to +2 cycles is not achievable at fixed times far from failure: on PHM08, eleven models reached at most 33.5\% at half of life. Letting the model decide when to speak changes the picture. A horizon-focused gradient-boosted regressor with a window-centred, debounced trigger, tuned by nested cross-validation, placed 91.7\% of PHM08 first warnings and 90.2\% of 927 simulated first warnings inside the window at a warning level of 20 cycles; under an identical protocol it outperformed re-implementations of the three winning PHM08 challenge approaches (significantly for two of them), and a Transformer given the same inputs did about as well. What this score measures is the alarm time: a constant estimate at the same times scored 90.8\%, whereas a classical health-indicator alarm scored 53.4\%. Priced in a renewal--reward model, window-tuned alarms were cheaper than optimal age replacement in 94 of 96 settings on the engines and batteries and had lower regret than alarms tuned directly to the cost, so the asymmetric window serves as a robust surrogate for the maintenance cost. Three pre-specified confirmations show where this holds: the method succeeded on measured battery degradation (97.6\%, all four hypotheses supported), fell short on N-CMAPSS (68.7\%), where the alarm time is less concentrated and the alarms were no cheaper than age replacement, and failed on rolling bearings (0.0\%). The approach is simple, fast and easy to audit; before it is relied on, the concentration of its alarm time at the intended warning level should be estimated on the fleet's own run-to-failure history. Section 8 outlines how the same alarm can set test and maintenance policies and serve automatic fault detection and diagnosis.

\section*{Declaration of generative AI and AI-assisted technologies in the manuscript preparation process}

\noindent During the preparation of this work the authors used Claude (Anthropic) in order to assist with code development, data analysis and drafting of the manuscript. After using this tool, the authors reviewed and edited the content as needed and take full responsibility for the content of the published article.

\section*{Data availability}

\noindent The PHM08, C-MAPSS and N-CMAPSS datasets are publicly available from the NASA Prognostics Center of Excellence data repository, the battery data of Severson et al. \cite{severson2019} from the data platform of the Toyota Research Institute (\url{https://data.matr.io}), and the bearing data of the IEEE PHM 2012 challenge \cite{nectoux2012} from the FEMTO-ST institute. Code, hashed protocols, derived data summaries and all result files needed to reproduce every number are deposited in a public repository; the link is withheld for anonymized review.

\begin{thebibliography}{34}
\bibitem{lei2018} Lei Y, Li N, Guo L, Li N, Yan T, Lin J. Machinery health prognostics: A systematic review from data acquisition to RUL prediction. Mechanical Systems and Signal Processing 2018;104:799--834. \url{https://doi.org/10.1016/j.ymssp.2017.11.016}
\bibitem{si2011} Si X-S, Wang W, Hu C-H, Zhou D-H. Remaining useful life estimation -- A review on the statistical data driven approaches. European Journal of Operational Research 2011;213(1):1--14. \url{https://doi.org/10.1016/j.ejor.2010.11.018}
\bibitem{ramasso2014} Ramasso E, Saxena A. Performance benchmarking and analysis of prognostic methods for CMAPSS datasets. International Journal of Prognostics and Health Management 2014;5(2). \url{https://doi.org/10.36001/ijphm.2014.v5i2.2236}
\bibitem{saxena2008} Saxena A, Goebel K, Simon D, Eklund N. Damage propagation modeling for aircraft engine run-to-failure simulation. In: 2008 International Conference on Prognostics and Health Management. IEEE; 2008, p. 1--9. \url{https://doi.org/10.1109/PHM.2008.4711414}
\bibitem{depater2022} de Pater I, Reijns A, Mitici M. Alarm-based predictive maintenance scheduling for aircraft engines with imperfect remaining useful life prognostics. Reliability Engineering \& System Safety 2022;221:108341. \url{https://doi.org/10.1016/j.ress.2022.108341}
\bibitem{saxena2010} Saxena A, Celaya J, Saha B, Saha S, Goebel K. Metrics for offline evaluation of prognostic performance. International Journal of Prognostics and Health Management 2010;1(1). \url{https://doi.org/10.36001/ijphm.2010.v1i1.1336}
\bibitem{severson2019} Severson KA, Attia PM, Jin N, Perkins N, Jiang B, Yang Z, Chen MH, Aykol M, Herring PK, Fraggedakis D, Bazant MZ, Harris SJ, Chueh WC, Braatz RD. Data-driven prediction of battery cycle life before capacity degradation. Nature Energy 2019;4(5):383--391. \url{https://doi.org/10.1038/s41560-019-0356-8}
\bibitem{ariaschao2021} Arias Chao M, Kulkarni C, Goebel K, Fink O. Aircraft engine run-to-failure dataset under real flight conditions for prognostics and diagnostics. Data 2021;6(1):5. \url{https://doi.org/10.3390/data6010005}
\bibitem{nectoux2012} Nectoux P, Gouriveau R, Medjaher K, Ramasso E, Chebel-Morello B, Zerhouni N, Varnier C. PRONOSTIA: An experimental platform for bearings accelerated degradation tests. In: IEEE International Conference on Prognostics and Health Management (PHM 2012). IEEE; 2012, p. 1--8. \url{https://hal.science/hal-00719503}
\bibitem{wang2008} Wang T, Yu J, Siegel D, Lee J. A similarity-based prognostics approach for remaining useful life estimation of engineered systems. In: 2008 International Conference on Prognostics and Health Management. IEEE; 2008, p. 1--6. \url{https://doi.org/10.1109/PHM.2008.4711421}
\bibitem{heimes2008} Heimes FO. Recurrent neural networks for remaining useful life estimation. In: 2008 International Conference on Prognostics and Health Management. IEEE; 2008, p. 1--6. \url{https://doi.org/10.1109/PHM.2008.4711422}
\bibitem{peel2008} Peel L. Data driven prognostics using a Kalman filter ensemble of neural network models. In: 2008 International Conference on Prognostics and Health Management. IEEE; 2008, p. 1--6. \url{https://doi.org/10.1109/PHM.2008.4711423}
\bibitem{zheng2017} Zheng S, Ristovski K, Farahat A, Gupta C. Long short-term memory network for remaining useful life estimation. In: 2017 IEEE International Conference on Prognostics and Health Management (ICPHM). IEEE; 2017, p. 88--95. \url{https://doi.org/10.1109/ICPHM.2017.7998311}
\bibitem{li2018} Li X, Ding Q, Sun J-Q. Remaining useful life estimation in prognostics using deep convolution neural networks. Reliability Engineering \& System Safety 2018;172:1--11. \url{https://doi.org/10.1016/j.ress.2017.11.021}
\bibitem{grinsztajn2022} Grinsztajn L, Oyallon E, Varoquaux G. Why do tree-based models still outperform deep learning on typical tabular data? Advances in Neural Information Processing Systems 2022;35:507--520. \url{https://doi.org/10.52202/068431-0037}
\bibitem{friedman2001} Friedman JH. Greedy function approximation: A gradient boosting machine. The Annals of Statistics 2001;29(5):1189--1232. \url{https://doi.org/10.1214/aos/1013203451}
\bibitem{rengasamy2020} Rengasamy D, Rothwell BC, Figueredo GP. Asymmetric loss functions for deep learning early predictions of remaining useful life in aerospace gas turbine engines. In: 2020 International Joint Conference on Neural Networks (IJCNN). IEEE; 2020, p. 1--7. \url{https://doi.org/10.1109/IJCNN48605.2020.9207051}
\bibitem{nguyen2019} Nguyen KTP, Medjaher K. A new dynamic predictive maintenance framework using deep learning for failure prognostics. Reliability Engineering \& System Safety 2019;188:251--262. \url{https://doi.org/10.1016/j.ress.2019.03.018}
\bibitem{mitici2023} Mitici M, de Pater I, Barros A, Zeng Z. Dynamic predictive maintenance for multiple components using data-driven probabilistic RUL prognostics: The case of turbofan engines. Reliability Engineering \& System Safety 2023;234:109199. \url{https://doi.org/10.1016/j.ress.2023.109199}
\bibitem{lee2023} Lee J, Mitici M. Deep reinforcement learning for predictive aircraft maintenance using probabilistic remaining-useful-life prognostics. Reliability Engineering \& System Safety 2023;230:108908. \url{https://doi.org/10.1016/j.ress.2022.108908}
\bibitem{kamariotis2024} Kamariotis A, Tatsis K, Chatzi E, Goebel K, Straub D. A metric for assessing and optimizing data-driven prognostic algorithms for predictive maintenance. Reliability Engineering \& System Safety 2024;242:109723. \url{https://doi.org/10.1016/j.ress.2023.109723}
\bibitem{gebraeel2005} Gebraeel NZ, Lawley MA, Li R, Ryan JK. Residual-life distributions from component degradation signals: A Bayesian approach. IIE Transactions 2005;37(6):543--557. \url{https://doi.org/10.1080/07408170590929018}
\bibitem{barlow1960} Barlow R, Hunter L. Optimum preventive maintenance policies. Operations Research 1960;8(1):90--100. \url{https://doi.org/10.1287/opre.8.1.90}
\bibitem{varma2006} Varma S, Simon R. Bias in error estimation when using cross-validation for model selection. BMC Bioinformatics 2006;7:91. \url{https://doi.org/10.1186/1471-2105-7-91}
\bibitem{kapoor2023} Kapoor S, Narayanan A. Leakage and the reproducibility crisis in machine-learning-based science. Patterns 2023;4(9):100804. \url{https://doi.org/10.1016/j.patter.2023.100804}
\bibitem{nosek2018} Nosek BA, Ebersole CR, DeHaven AC, Mellor DT. The preregistration revolution. Proceedings of the National Academy of Sciences 2018;115(11):2600--2606. \url{https://doi.org/10.1073/pnas.1708274114}
\bibitem{pedregosa2011} Pedregosa F, Varoquaux G, Gramfort A, Michel V, Thirion B, Grisel O, Blondel M, Prettenhofer P, Weiss R, Dubourg V, VanderPlas J, Passos A, Cournapeau D, Brucher M, Perrot M, Duchesnay E. Scikit-learn: Machine learning in Python. Journal of Machine Learning Research 2011;12:2825--2830.
\bibitem{clopper1934} Clopper CJ, Pearson ES. The use of confidence or fiducial limits illustrated in the case of the binomial. Biometrika 1934;26(4):404--413. \url{https://doi.org/10.1093/biomet/26.4.404}
\bibitem{mcnemar1947} McNemar Q. Note on the sampling error of the difference between correlated proportions or percentages. Psychometrika 1947;12(2):153--157. \url{https://doi.org/10.1007/BF02295996}
\bibitem{holm1979} Holm S. A simple sequentially rejective multiple test procedure. Scandinavian Journal of Statistics 1979;6(2):65--70. \url{https://www.jstor.org/stable/4615733}
\bibitem{dejonge2020} de Jonge B, Scarf PA. A review on maintenance optimization. European Journal of Operational Research 2020;285(3):805--824. \url{https://doi.org/10.1016/j.ejor.2019.09.047}
\bibitem{venkatasubramanian2003} Venkatasubramanian V, Rengaswamy R, Kavuri SN, Yin K. A review of process fault detection and diagnosis: Part III: Process history based methods. Computers \& Chemical Engineering 2003;27(3):327--346. \url{https://doi.org/10.1016/S0098-1354(02)00162-X}
\bibitem{page1954} Page ES. Continuous inspection schemes. Biometrika 1954;41(1--2):100--115. \url{https://doi.org/10.1093/biomet/41.1-2.100}
\bibitem{roberts1959} Roberts SW. Control chart tests based on geometric moving averages. Technometrics 1959;1(3):239--250. \url{https://doi.org/10.2307/1271439}
\end{thebibliography}
\end{document}
